\section{Related Work}
\label{sec:related-work}

Memory management in \ac{HPC} environments has been extensively studied, with research spanning predictive modeling, runtime adaptation, and distributed scheduling strategies.
This section reviews key approaches and highlights how our work advances the state of the art.

\subsection{Memory Prediction in HPC}

Several studies have addressed memory prediction for \ac{HPC} workloads.
Rodrigues et al.~\cite{rodrigues2016} employed machine learning to classify job memory usage into bins based on features such as user ID, number of processors, and job names, achieving up to 98\% accuracy.
However, their approach relies on historical scheduler data and operates at job-level granularity, limiting applicability to new users or novel workflows.

Li et al.~\cite{li2019} developed a two-stage prediction model that first classifies jobs as high-memory or normal, then applies specialized regressors for extreme cases.
While effective for outlier detection, this approach still depends on job logs rather than data-level features, preventing runtime adjustments for data-intensive tasks.

\subsection{Memory-Aware Scheduling}

Memory-aware scheduling has gained attention in distributed computing frameworks.
Khandelwal et al.~\cite{khandelwal2020} presented Cleo, which profiles application memory usage and migrates memory-intensive containers.
Thamsen et al.~\cite{thamsen2017} introduced Mary, a system that monitors resource usage and adjusts allocations dynamically.
These systems focus on container-level management rather than data partitioning strategies.

For data-parallel frameworks specifically, Shi et al.~\cite{shi2014} proposed MRTuner for automatic MapReduce configuration, including memory settings.
However, their approach targets Hadoop environments and does not address the unique challenges of scientific computing workloads with complex memory patterns.

\subsection{Chunking Strategies in Distributed Computing}

Chunking strategies have been explored primarily in the context of \ac{I/O} optimization and cache efficiency.
Zhang et al.~\cite{zhang2019} developed adaptive chunking for scientific datasets stored in HDF5, focusing on storage layout rather than runtime memory constraints.
Tantisiriroj et al.~\cite{tantisiriroj2011} examined chunking for parallel file systems but did not consider operator-specific memory requirements.

Dask's auto-chunking mechanism~\cite{rocklin2015dask} represents the current state of practice for Python-based distributed computing.
It employs simple heuristics based on target chunk sizes (typically 128MB) without considering algorithm complexity or intermediate memory expansion.
This limitation is particularly problematic for tensor operations and scientific computing workflows where memory usage can vary dramatically based on operator characteristics.

\subsection{Performance Modeling in Scientific Computing}

Performance modeling for scientific applications has been extensively studied.
Carrington et al.~\cite{carrington2006} developed performance models for \ac{HPC} applications using machine learning, though focusing on execution time rather than memory.
Bhimji et al.~\cite{bhimji2016} analyzed \ac{I/O} patterns in scientific workflows but did not address memory constraints in distributed settings.

For Python-specific workloads, Harris et al.~\cite{harris2020} demonstrated that array computing libraries exhibit predictable performance patterns.
However, their work focused on single-node execution rather than distributed memory management.

\subsection{Memory Measurement Challenges}

Accurate memory measurement in Python presents unique challenges due to:
\begin{itemize}
    \item Reference counting and garbage collection overhead
    \item C extension memory allocations invisible to Python
    \item Operating system page caching and memory sharing
    \item Container virtualization layers
\end{itemize}

Our work addresses these challenges through comprehensive profiling infrastructure that captures memory usage at multiple system levels, ensuring accurate predictions for distributed execution.

\subsection{Distinguishing Features of Our Approach}

Our work differs from existing approaches in several key aspects:

\begin{enumerate}
    \item \textbf{Shape-based prediction}: We model memory consumption based on input tensor dimensions and operator characteristics, enabling accurate predictions for scientific workloads
    \item \textbf{Integration with existing frameworks}: Our solution extends Dask's chunking mechanism rather than requiring a new runtime system
    \item \textbf{Operator-aware modeling}: We account for algorithm-specific memory patterns, including intermediate expansions common in tensor operations
    \item \textbf{Practical validation}: We demonstrate effectiveness on real seismic processing workloads with complex memory requirements
\end{enumerate}

Table~\ref{tab:comparison} summarizes how our approach compares to representative prior work.

\begin{table}[ht]
\centering
\caption{Comparison with Related Work}
\label{tab:comparison}
\begin{tabular}{lcccc}
\toprule
\textbf{Approach} & \textbf{Prediction} & \textbf{Data-} & \textbf{Operator-} & \textbf{Runtime} \\
                  & \textbf{Basis}      & \textbf{aware} & \textbf{specific}  & \textbf{Adjust.} \\
\midrule
Rodrigues~\cite{rodrigues2016} & Job logs    & No  & No  & No  \\
Li~\cite{li2019}               & Job logs    & No  & No  & No  \\
Khandelwal~\cite{khandelwal2020} & Runtime  & No  & No  & Yes \\
Dask auto~\cite{rocklin2015dask}  & File size & Yes & No  & No  \\
\textbf{Our work}              & Shape+Op    & Yes & Yes & Yes \\
\bottomrule
\end{tabular}
\end{table}